% Contribution Template — NCA6 Coauthor Section
%
% HOW TO USE:
% 1. Copy this file: cp _template.tex [your-topic].tex
%    (e.g., labor_supply.tex, health_welfare.tex, trade_flows.tex)
% 2. Replace every [BRACKETED] placeholder with your content.
% 3. Add \input{sections/contributions/[your-topic].tex} to paper/main.tex.
% 4. Verify your KMs all carry confidence labels (and likelihood when probabilistic).
% 5. Run /review on this file when you're ready for critique.
%
% Format follows NCA convention: Section, KMs (each with Evidence Base,
% Uncertainties, Confidence, Likelihood), Traceable Account.

\contribution{[Section Title — e.g., Climate Impacts on Labor Supply]}{[Your Name]}

% ============================================================
% Opening — 1–2 paragraphs setting up the section
% ============================================================
[Open with a short framing of your slice of climate-economics. What is in
scope, what is out. Connect to one or two of the anchors (NCA5 Ch.19, IPCC
AR6 chapters) and indicate where your contribution extends, refines, or
contests prior findings.]

% ============================================================
% Key Messages
% ============================================================

\begin{keymessage}{[KM-1 short title]}
[One paragraph stating the key message in plain language. Avoid jargon where
possible; readers come from many subfields.] \confidence{[very high / high / medium / low]}
\likelihood{[very likely / likely / as likely as not / unlikely / very unlikely]}
\end{keymessage}

\evidencebase
[2–4 paragraphs describing the empirical evidence supporting this KM. Cite
extensively from your bridge\_review.md output. Distinguish between
peer-reviewed papers, working papers, and meta-analyses.]

\uncertainties
[1–2 paragraphs on what remains uncertain. Be explicit: data limitations,
identification challenges, scope of generalization, model dependencies.]

\confidencedesc
[1 paragraph explaining \emph{why} you assign the confidence level you do.
Reference the strength of evidence and degree of consensus criteria.]

\likelihooddesc
[1 paragraph (if the claim is probabilistic) explaining the likelihood
assignment based on the supporting literature.]

\traceableaccount
[Bullet list connecting each clause of the KM to specific citations. This
is the auditable evidence chain — a future reader should be able to trace
every claim back to its source.]
\begin{itemize}
  \item Claim X $\rightarrow$ \cite{[citekey1]}, \cite{[citekey2]}
  \item Claim Y $\rightarrow$ \cite{[citekey3]}
  \item Magnitude estimate $\rightarrow$ \cite{[citekey4]} (point estimate),
        \cite{[citekey5]} (range)
\end{itemize}

% ----- Repeat for KM-2, KM-3, ... -----

\begin{keymessage}{[KM-2 short title]}
[...] \confidence{[level]}
\end{keymessage}

\evidencebase
[...]

\uncertainties
[...]

\confidencedesc
[...]

\traceableaccount
\begin{itemize}
  \item [...]
\end{itemize}

% ============================================================
% Optional — section-level synthesis
% ============================================================

\subsection*{Synthesis and Open Questions}

[Optional closing — pulls together the KMs in this contribution into a
short narrative and notes what your contribution leaves open for future
NCA cycles or for adjacent coauthor contributions to address.]
